\section{\texorpdfstring{Равномерная сходимость. Пространство \(\ell^\infty(\Omega)\)}{Равномерная сходимость. Пространство ограниченных функций}}\label{sec-sequences-01}

Напоминание: нормированное пространство, полное относительно метрики, порождённой нормой, называется \textbf{банаховым}.

В конечномерном пространстве эквивалентны все \emph{нормы}. В бесконечномерном пространстве нормы могут быть неэквивалентны.

\begin{definition}[Пространство ограниченных функций]\label{def-bounded}

\(\ell^\infty(\Omega)\) --- множество функций \(f:\Omega\to\mathbb C\), ограниченных по модулю: \[
\exists C>0\quad\forall x\in\Omega:\quad |f(x)|\le C.
\] На нём задаётся норма

\begin{equation}
\lVert f\rVert_\infty=\sup_{x\in\Omega}|f(x)|,
\qquad d(f,g)=\lVert f-g\rVert_\infty.
\label{eq-sup-norm}
\end{equation}

Метрика \(d\) называется чебышевским расстоянием.

\end{definition}

\begin{definition}[Равномерная сходимость]\label{def-uniform-sequence}

Последовательность \(\{f_n\}_{n=1}^{\infty}\) равномерно сходится на \(\Omega\) к \(f\), если \[
\lVert f_n-f\rVert_\infty\longrightarrow0.
\] Эквивалентная запись:

\begin{equation}
\forall\varepsilon>0\ \exists N\ \forall n>N\ \forall x\in\Omega:
\quad |f_n(x)-f(x)|<\varepsilon.
\label{eq-uniform-quantifiers}
\end{equation}

Обозначение: \(f_n\rightrightarrows f\) на \(\Omega\).

\end{definition}

\begin{definition}[Поточечная сходимость]\label{def-pointwise}

Для фиксированного \(x\) числовая последовательность \(f_n(x)\) может сходиться к \(f(x)\). Множество таких точек называется множеством поточечной сходимости. На нём \[
\forall x\ \forall\varepsilon>0\ \exists N_x\ \forall n>N_x:
\quad |f_n(x)-f(x)|<\varepsilon.
\] В отличие от условия \eqref{eq-uniform-quantifiers}, номер \(N_x\) здесь может зависеть от точки.

\end{definition}

\begin{theorem}[Полнота \(\ell^\infty(\Omega)\)]\label{thm-completeness}

Пространство \(\ell^\infty(\Omega)\) полно по норме \(\lVert \cdot\rVert_\infty\).

\end{theorem}

\begin{proof}

Пусть \(\{f_n\}\) фундаментальна: \[
\forall\varepsilon>0\ \exists N\ \forall n,k>N:
\quad \lVert f_n-f_k\rVert_\infty<\varepsilon.
\] Для каждого \(x\) последовательность \(\{f_n(x)\}\) фундаментальна в \(\mathbb C\). По полноте \(\mathbb C\) существует \(f(x)=\lim_n f_n(x)\).

Сначала покажем ограниченность норм \(\lVert f_n\rVert_\infty\). В условии фундаментальности возьмём \(\varepsilon=1\) и положим \[
M=\max\{\lVert f_1\rVert_\infty,\ldots,\lVert f_{N+1}\rVert_\infty\}.
\] Для \(n>N\) имеем \(\lVert f_n\rVert_\infty\le\lVert f_n-f_{N+1}\rVert_\infty+\lVert f_{N+1}\rVert_\infty\le1+M\). Поэтому все \(f_n\) ограничены одной константой; переход к поточечному пределу даёт \(|f(x)|\le1+M\). Значит, \(f\in\ell^\infty(\Omega)\).

Теперь для заданного \(\varepsilon>0\) выбираем \(N\) с \(\lVert f_n-f_m\rVert_\infty<\varepsilon/2\) при \(m,n>N\). Переходя к пределу \(m\to\infty\) для каждого \(x\), получаем \[
|f_n(x)-f(x)|\le\varepsilon/2,
\qquad \lVert f_n-f\rVert_\infty\le\varepsilon/2<\varepsilon.
\] Итак, \(f_n\rightrightarrows f\).

\end{proof}
